
@@M0@@
Rappelons la définition du degré @@M1@@ de @@M2@@ [3] : il est nul sauf dans le cas où @@M3@@. Dans ce dernier cas on pose
@@M4@@
où @@M5@@ est le point générique de @@M6@@ et @@M7@@ le point générique de @@M8@@.
\end{enumerate}

Dans la démonstration de (4.4.3) le cas @@M9@@ est le seul non-trivial. Si @@M10@@, alors on a
@@M11@@
pour @@M12@@ dans un voisinage de @@M13@@. Mais la condition (4.4.5) est satisfaite aussi dans le cas @@M14@@, car alors la dimension de la fibre générique est nulle (EGA III 4.2.2 et EGA IV 5.6.6). Alors 1.1.2 donne les congruences
@@M15@@
modulo @@M16@@. En vertu de (4.4.2) il en découle que
@@M17@@
D'autre part @@M18@@ donne
@@M19@@
d'où l'assertion.

